% DRAFT ONLY: original source visually checked; root audit pending.
\begin{saresult}{A.126}{Known}{Simply connected closed five-manifolds have trivial concordance components}
Let $M$ be a closed connected PL five-manifold with $\pi_1M=0$.
For any definable realization of its PL model over any real closed
field,
\[
 \pi_0P_{\rm PL}(M)=\pi_0P_{\rm def}(M)=0.
\]
There is no restriction on $\pi_2M$. The smooth input holds for
closed simply connected smooth manifolds in every dimension at
least five. Only the closed case of that input is invoked here.
\end{saresult}
\begin{proof}
The needed theorem is the full pseudo-isotopy theorem, not merely
the implication that pseudo-isotopic endpoint diffeomorphisms are
isotopic. In the original paper, Cerf defines the pseudo-isotopy
group on printed p.~7 as the diffeomorphisms of $V\times I$
fixed on $V\times\{0\}$. Theorem~0 on printed p.~8 states that
this group is connected for closed smooth $V$ with
$\pi_1V=0$ and $\dim V\geq5$
\cite[Introduction, Definition~2 and Theorem~0, pp.~7--8]{Cerf1970}.
The endpoint assertion occurs separately as Corollary~1 on that
page. The main connectivity theorem is proved as Theorem~3 on
printed p.~156, using the critical-point-free function-space
model described in the introduction.

The smooth group with the bottom fixed is locally path connected,
so this connectedness means that its component group is zero.
For completeness, choose a product metric near the boundary of
$V\times I$. Sufficiently $C^1$-small diffeomorphisms are expressed
uniquely as $x\mapsto\exp_x X(x)$ in the usual local-addition
chart. The boundary-preserving condition gives vector fields
tangent to the boundary, and the bottom-fixed condition gives
$X=0$ on the bottom. A small convex neighborhood of zero in this
linear space stays inside the diffeomorphism chart. Scaling $X$
to zero therefore gives paths through bottom-fixed smooth
diffeomorphisms. This proves the required local path connectedness
in the $C^\infty$ topology. Consequently Cerf's theorem gives
\[
 \pi_0P_{\rm Diff}(V)=0.                              \tag{A.126.1}
\]

Every PL five-manifold admits a compatible smoothing, without any
independently prescribed boundary smoothing. Here is the exact
low-dimensional input already used in A.124. In a finite PL handle
decomposition the index-$k$ extension obstruction is
$\Gamma_{k-1}$
\cite[Essay~IV, Appendix~B, p.~211]{KirbySiebenmann1977}.
The compatible low-dimensional smoothings in
\cite[Essay~I, p.~39]{KirbySiebenmann1977}, together with
$\Gamma_4=0$ explicitly on
\cite[Essay~IV, p.~200]{KirbySiebenmann1977}, remove all obstructions
for handles of index at most five. Let $M^{\rm sm}$ be a compatible
smoothing of $M$ and $\tau:|K|\to M^{\rm sm}$ its differentiable
triangulation. It is still closed and simply connected.

The precise category comparison now used is the component
isomorphism of Burghelea--Lashof
\cite[Theorem~6.1, p.~41, and proof pp.~42--43]{BurgheleaLashof1974II},
with the concordance groups defined on p.~41:
\[
 c_\tau:\pi_0P_{\rm Diff}(M^{\rm sm})
       \xrightarrow{\ \cong\ }\pi_0P_{\rm PL}(|K|).    \tag{A.126.2}
\]
Their bottom-and-side convention agrees with ours; the side is
empty here. This is the unstable component comparison associated
to the differentiable triangulation, as already used in A.103 and
A.124. It does not require a smooth map to be literally PL in a
fixed coordinate triangulation, and no stable comparison is being
substituted for (A.126.2). Equations (A.126.1)--(A.126.2) give
$\pi_0P_{\rm PL}(M)=0$.

Finally the finite PL field comparison A.115 and the relative
definable component comparison A.113 preserve this group and its
bottom condition. They prove the asserted conclusion over every
real closed field. The absence of a $\pi_2$ assumption is exactly
what extends the dimension-five range of A.124.
\end{proof}
